\documentclass[11pt]{article}

\usepackage[T1]{fontenc}
\usepackage[utf8]{inputenc}
\usepackage{lmodern}
\usepackage{microtype}
\usepackage[margin=1in]{geometry}
\usepackage{amsmath,amssymb,amsthm,mathtools}
\usepackage{enumitem}
\usepackage{hyperref}

\hypersetup{
  colorlinks=true,
  linkcolor=blue,
  urlcolor=blue,
  citecolor=blue
}

\newtheorem{theorem}{Theorem}
\newtheorem{lemma}[theorem]{Lemma}
\newtheorem{corollary}[theorem]{Corollary}
\newtheorem{proposition}[theorem]{Proposition}
\theoremstyle{definition}
\newtheorem{definition}[theorem]{Definition}
\theoremstyle{remark}
\newtheorem{remark}[theorem]{Remark}

\newcommand{\F}{\mathbb{F}_2}
\newcommand{\Pn}{\mathcal{P}_n}
\newcommand{\Amb}{\operatorname{Amb}}
\newcommand{\ip}[2]{\langle #1,#2\rangle}
\newcommand{\ind}{\mathbf{1}}

\title{Efficient Online Learning of Pairwise-Ambiguous Parities\\
\large and an LPN-Based Computational Separation}
\author{}
\date{}

\begin{document}
\maketitle

\begin{abstract}
This note explains an efficient online algorithm for a structured ambiguous
version of the parity class.  The unknown target consists of two parity
functions.  On every instance, either parity output is acceptable, and the
adversary may reveal either acceptable label.  Even though the learner is not
told whether its own prediction was acceptable, a quadratic feature-map
argument gives a deterministic polynomial-time learner making at most
$\binom{n+2}{2}-1$ mistakes.

The result should not be interpreted as an efficient algorithm for ordinary
Learning Parity with Noise (LPN).  Rather, LPN is used as a hardness
assumption: pairwise-ambiguous parities are efficiently learnable in the
mistake-unknown model, while efficient agnostic online learning of the
ordinary parity class would yield a polynomial-time distinguisher for
decisional LPN.
\end{abstract}

\section{The distinction from ordinary LPN}

Let
\[
\Pn
=
\left\{
  p_w(x)=\ip{w}{x}\bmod 2 : w\in\F^n
\right\}
\]
be the class of homogeneous parity functions on $\F^n$.

In ordinary LPN, examples have the form
\[
  y=p_s(x)\oplus e,
  \qquad
  e\sim\operatorname{Bernoulli}(\eta),
\]
where the noise bit is random and independent of $x$.  The efficient result
below concerns a different model.  There are two fixed latent parities
$p_u$ and $p_v$, and the acceptable-label set on $x$ is
\[
  \{p_u(x),p_v(x)\}.
\]
Thus, relative to $p_u$, a revealed label can be written as
\[
  y=p_u(x)\oplus a\,\ip{u+v}{x},
  \qquad a\in\{0,1\},
\]
where the adversary may choose $a$ separately on every round.  A flip is
possible only on the structured disagreement set
\[
  \{x:\ip{u+v}{x}=1\}.
\]
This is not independent random classification noise.

\section{Pairwise ambiguity of parities}

\begin{definition}[Pairwise-ambiguous parity class]
For $u,v\in\F^n$, define the set-valued hypothesis
\[
  G_{u,v}(x)=\{p_u(x),p_v(x)\}.
\]
The pairwise-ambiguous closure of the parity class is
\[
  \Amb^{(2)}(\Pn)
  =
  \{G_{u,v}:u,v\in\F^n\}.
\]
\end{definition}

At points where $p_u(x)=p_v(x)$, the acceptable-label set is a singleton.
At points where they disagree, both labels are acceptable:
\[
G_{u,v}(x)=
\begin{cases}
\{p_u(x)\}, & p_u(x)=p_v(x),\\
\{0,1\}, & p_u(x)\neq p_v(x).
\end{cases}
\]
The disagreement region is
\[
A_{u,v}
=
\{x:p_u(x)\neq p_v(x)\}
=
\{x:\ip{u+v}{x}=1\}.
\]
When $u\neq v$, this is an affine hyperplane of $\F^n$; when $u=v$, it is
empty.  Hence the ambiguity is highly structured rather than arbitrary.

\section{The mistake-unknown online protocol}

Fix an unknown $G_{u,v}\in\Amb^{(2)}(\Pn)$.  On round $t$:
\begin{enumerate}[label=\arabic*.]
  \item The adversary presents an instance $x_t\in\F^n$.
  \item The learner predicts $\widehat y_t\in\{0,1\}$.
  \item The adversary reveals a label
  \[
    y_t\in G_{u,v}(x_t).
  \]
\end{enumerate}
The learner makes a mistake exactly when
\[
  \widehat y_t\notin G_{u,v}(x_t).
\]
It observes the valid demonstration $(x_t,y_t)$, but receives no bit telling
it whether $\widehat y_t$ was valid.  In particular,
$\widehat y_t\neq y_t$ does not imply a mistake: both labels might have been
acceptable.

\section{A quadratic representation of validity}

All arithmetic in this section is over $\F$.  Define
\[
q_{u,v}(x,y)
=
\bigl(y+\ip{u}{x}\bigr)
\bigl(y+\ip{v}{x}\bigr).
\]

\begin{lemma}[Validity polynomial]
For every $x\in\F^n$ and $y\in\F$,
\[
  y\in G_{u,v}(x)
  \quad\Longleftrightarrow\quad
  q_{u,v}(x,y)=0.
\]
\end{lemma}

\begin{proof}
The product is zero exactly when at least one factor is zero.  This happens
exactly when $y=\ip{u}{x}$ or $y=\ip{v}{x}$.
\end{proof}

On Boolean inputs, $x_i^2=x_i$ and $y^2=y$.  After multilinearizing the
quadratic polynomial, define the feature map
\[
\phi(x,y)
=
\left(
  y,
  (x_i)_{i=1}^n,
  (yx_i)_{i=1}^n,
  (x_ix_j)_{1\leq i<j\leq n}
\right)
\in\F^D,
\]
where
\[
D
=1+2n+\binom{n}{2}
=\binom{n+2}{2}.
\]
There is a coefficient vector $\theta_{u,v}\in\F^D$ such that
\[
  q_{u,v}(x,y)=\ip{\theta_{u,v}}{\phi(x,y)}.
\]
Moreover, $\theta_{u,v}\neq 0$: the coefficient of the feature $y$ is always
one, coming from $y^2=y$.

Consequently, the set of valid labeled examples is contained in the
codimension-one subspace
\[
  \theta_{u,v}^{\perp}
  =
  \{z\in\F^D:\ip{\theta_{u,v}}{z}=0\}.
\]

\section{The span learner}

Before round $t$, maintain
\[
R_t
=
\operatorname{span}
\{\phi(x_s,y_s):s<t\}
\subseteq\F^D.
\]
The learner acts as follows.

\begin{enumerate}[label=\arabic*.]
  \item If $\phi(x_t,0)\in R_t$, predict $\widehat y_t=0$.
  \item Otherwise, if $\phi(x_t,1)\in R_t$, predict
        $\widehat y_t=1$.
  \item Otherwise, predict either label, say $\widehat y_t=0$.
  \item After observing the valid label $y_t$, update
  \[
    R_{t+1}
    =
    \operatorname{span}\bigl(R_t\cup\{\phi(x_t,y_t)\}\bigr).
  \]
\end{enumerate}

\begin{lemma}[Soundness of span certificates]
If $\phi(x,b)\in R_t$, then $b\in G_{u,v}(x)$.
\end{lemma}

\begin{proof}
Every revealed label is valid, so for every $s<t$,
\[
  \ip{\theta_{u,v}}{\phi(x_s,y_s)}=0.
\]
Therefore $R_t\subseteq\theta_{u,v}^{\perp}$.  If
$\phi(x,b)\in R_t$, then
\[
  q_{u,v}(x,b)
  =
  \ip{\theta_{u,v}}{\phi(x,b)}
  =0.
\]
The validity-polynomial lemma now gives $b\in G_{u,v}(x)$.
\end{proof}

\section{Mistake bound and running time}

\begin{theorem}[Efficient learning of pairwise-ambiguous parities]
The span learner is deterministic, runs in time polynomial in $n$ per round,
and makes at most
\[
  D-1
  =
  \binom{n+2}{2}-1
  =
  \frac{n(n+3)}{2}
\]
mistakes on every sequence realizable by a fixed
$G_{u,v}\in\Amb^{(2)}(\Pn)$.
\end{theorem}

\begin{proof}
First, every maintained vector is orthogonal to $\theta_{u,v}$, so
\[
  R_t\subseteq\theta_{u,v}^{\perp}.
\]
Since $\theta_{u,v}\neq 0$,
\[
  \dim(\theta_{u,v}^{\perp})=D-1.
\]

Suppose the learner makes a mistake on round $t$.  We claim that the newly
revealed valid vector $\phi(x_t,y_t)$ was not already in $R_t$.  If it were in
$R_t$, then one of two things would happen.  Either the learner would predict
$y_t$, or the learner would predict the other label because that label also
had its feature vector in $R_t$.  In the latter case, the soundness lemma says
the other label is valid as well.  In both cases the prediction would be
acceptable, contradicting that a mistake occurred.

Thus every mistake satisfies
\[
  \phi(x_t,y_t)\notin R_t,
\]
and therefore causes a strict dimension increase:
\[
  \dim(R_{t+1})=\dim(R_t)+1.
\]
The dimension can increase at most $D-1$ times while the span remains inside
$\theta_{u,v}^{\perp}$.  Hence the total number of mistakes is at most $D-1$.

For computation, maintain a row-echelon basis of $R_t$ over $\F$.  The feature
dimension is $D=\Theta(n^2)$, and membership tests and insertion can be done
by incremental Gaussian elimination.  A direct implementation uses
$O(D^2)=O(n^4)$ bit operations per round, which is polynomial in $n$.
\end{proof}

\begin{remark}[Information-theoretic versus efficient bounds]
There are at most $2^{2n}$ ordered pairs $(u,v)$.  A generic finite-class
argument for the mistake-unknown model can therefore give an abstract
$2n$-type mistake bound.  That argument does not by itself provide an
efficient implicit implementation over the exponentially large class.  The
quadratic-span learner gives an explicit polynomial-time algorithm, at the
cost of the weaker $O(n^2)$ mistake bound.
\end{remark}

\section{The LPN-based computational separation}

We now explain the role of LPN.  Fix a constant noise rate
$\eta\in(0,1/2)$.  The decisional-LPN problem is to distinguish the following
two distributions over labeled examples:
\begin{align*}
\mathsf{LPN}_{s,\eta}:&
&x&\sim\operatorname{Unif}(\F^n),
&y&=\ip{s}{x}\oplus e,
&e&\sim\operatorname{Bernoulli}(\eta),\\
\mathsf{Unif}:&
&x&\sim\operatorname{Unif}(\F^n),
&y&\sim\operatorname{Unif}(\F),
\end{align*}
where $s\in\F^n$ is an unknown secret.

\begin{proposition}[Agnostic online parity learning would break LPN]
Suppose that, for some constant $\alpha>0$ and polynomial $p$, there is a
randomized online learner for $\Pn$ that runs in time polynomial in $n$ and
$T$ and has expected regret
\[
  R_T\leq p(n)T^{1-\alpha}
\]
against the best parity on every labeled sequence.  Then decisional LPN at
any fixed noise rate $\eta<1/2$ can be solved in randomized polynomial time.
The learner may be improper.
\end{proposition}

\begin{proof}
Given $T$ labeled samples, feed them to the online learner in sequence and
let
\[
  L_A
  =
  \sum_{t=1}^T \ind\{\widehat y_t\neq y_t\}
\]
be its ordinary binary prediction loss.

Under $\mathsf{LPN}_{s,\eta}$, the comparator $p_s$ has expected loss
$\eta T$.  Therefore the regret guarantee implies
\[
  \mathbb{E}[L_A\mid\mathsf{LPN}_{s,\eta}]
  \leq
  \eta T+R_T.
\]
Under $\mathsf{Unif}$, each label is a fresh unbiased bit, independent of the
learner's prediction conditional on the past.  Hence
\[
  \mathbb{E}[L_A\mid\mathsf{Unif}]=\frac{T}{2}.
\]

Turn this expectation gap into a distinguisher as follows: after computing
$L_A$, output ``uniform'' with probability $L_A/T$, and output ``LPN''
otherwise.  Its probability of outputting ``uniform'' is exactly $1/2$ under
$\mathsf{Unif}$ and at most
\[
  \eta+\frac{R_T}{T}
\]
under $\mathsf{LPN}_{s,\eta}$.  Thus its distinguishing advantage is at least
\[
  \frac12-\eta-\frac{R_T}{T}.
\]
Choose a polynomially large horizon $T$ such that
\[
  \frac{R_T}{T}
  \leq
  \frac{1/2-\eta}{2}.
\]
This is possible because
$R_T/T\leq p(n)T^{-\alpha}$ and $\alpha$ is a fixed positive constant.  The
resulting advantage is a positive constant, and the total running time and
sample size are polynomial.
\end{proof}

\begin{corollary}[Conditional separation]
Under the standard decisional-LPN hardness assumption,
\begin{enumerate}[label=(\roman*)]
  \item $\Amb^{(2)}(\Pn)$ is polynomial-time learnable in the realizable
        mistake-unknown model, with at most $n(n+3)/2$ mistakes; but
  \item $\Pn$ has no polynomial-time agnostic online learner, proper or
        improper, with regret $p(n)T^{1-\alpha}$ for any fixed
        $\alpha>0$ and polynomial $p$.
\end{enumerate}
\end{corollary}

The first item is unconditional.  Only the hardness claim in the second item
is conditional on LPN.

\section{Why the ambiguous learner does not solve ordinary LPN}

In a fully arbitrary ambiguity model, one could declare an arbitrary set of
instances ambiguous and thereby hide an arbitrary corruption pattern.  That
would risk making ambiguity as hard as agnostic noise.

For pairwise parities, the ambiguity region is restricted to
\[
  A_d=\{x:\ip{d}{x}=1\}
\]
for one fixed $d=u+v$.  An independent LPN noise pattern is overwhelmingly
unlikely to agree with one such fixed parity-defined region across all
examples.  Therefore the polynomial-time learner for
$\Amb^{(2)}(\Pn)$ does not imply a polynomial-time algorithm for LPN.

The result instead identifies a structured form of adversarial uncertainty
that is tractable even though unrestricted agnostic parity learning is
LPN-hard.

\section{Extension to constantly many latent parities}

For a fixed integer $r\geq 2$, define
\[
  G_{w_1,\ldots,w_r}(x)
  =
  \{\ip{w_1}{x},\ldots,\ip{w_r}{x}\}.
\]
Validity is characterized by
\[
  q(x,y)
  =
  \prod_{j=1}^r\bigl(y+\ip{w_j}{x}\bigr),
  \qquad
  y\in G_{w_1,\ldots,w_r}(x)
  \Longleftrightarrow
  q(x,y)=0.
\]
After multilinearization, $q$ is represented using all nonconstant
square-free monomials of degree at most $r$ in the $n+1$ Boolean variables
$x_1,\ldots,x_n,y$.  The feature dimension is at most
\[
  D_r
  =
  \sum_{k=1}^r\binom{n+1}{k}
  =
  n^{O(r)}.
\]
Because all homogeneous parities vanish at $x=0$, we have $q(0,1)=1$, so the
coefficient vector is nonzero.  The same span argument therefore gives at
most $D_r-1$ mistakes and an $n^{O(r)}$-time implementation.  Consequently,
for every fixed constant $r$, the class of ambiguities generated by $r$
latent parities is efficiently mistake-unknown learnable.

\section{Relation to the subspace formulation}

For $r=2$, write $d=u+v$ and $U=\operatorname{span}\{d\}$.  Relative to the
base parity $p_u$, the two acceptable labels are obtained by adding either
zero or an element of the one-dimensional subspace $U$ evaluated at $x$.
Thus the pairwise construction is exactly the special case in which the
ambiguity subspace has dimension at most one.

\section{Conclusion}

The result is an efficient learning theorem for \emph{pairwise-ambiguous
parities}, not an algorithm for ordinary noisy parity learning.  Its core is
a low-degree algebraic representation: every valid labeled example lies in
one unknown codimension-one subspace of a quadratic feature space.  A learner
can accumulate revealed valid examples, use span membership as a sound
certificate, and charge every actual mistake to a dimension increase.

Together with the decisional-LPN assumption, this gives a clean computational
separation between structured mistake-unknown ambiguity and agnostic online
learning of the underlying parity class.

\end{document}
